% =====================================================================
% Introduction — contribution list. Scope-locked to sections (no claim
% stronger than the corresponding §). Mirrors frozen_claim_hierarchy_v1.md.
% =====================================================================
\paragraph{Contributions.}
\begin{itemize}
  \item \textbf{An output-interface view of language-grounded routing.} We frame constraint grounding as a
        choice of \emph{output representation} and show, on a constructed method-agnostic held-out
        closed-loop SUMO stress test (SUMO-Disc-Stress-v1, cases selected by graph geometry and oracle
        constraints only), that a dense $\Theta(|E|)$ edge-weight interface collapses into a
        fixed-capacity, scenario-invariant enumeration (100\% single near-uniform outputs), failing both
        constraint satisfaction and efficiency.
  \item \textbf{STAG-Route: typed sparse anchors.} Grounding constraints into typed sparse anchors
        preserves edge-level selectivity and route capacity, attaining conditional constraint satisfaction
        $0.96$ at low excess detour and ranking first of nine methods on satisfaction and travel time.
  \item \textbf{A structural separation from numeric sensing.} A near-oracle numeric GNN baseline fails on
        the multi-edge buckets that define the task, and a coverage ablation shows this failure is
        \emph{invariant} to sensor fidelity---isolating constraint \emph{representation} from sensing.
  \item \textbf{Mechanism and robustness ablations.} We localize the grounding to specific semantic fields
        (direction for localization; event-type/severity for conflict resolution), show the verifier
        normalizes rather than repairs an already hallucination-free model, and show the typed-anchor
        interface is robust across model families and across held-out topologies, with the closed-loop
        advantage transferring to unseen networks ($119/120$ vs.\ $0/120$ dense); segment-/edge-exact
        precision remains open.
  \item \textbf{An honest frontier.} We characterize and reconcile the method's limit: grounding is
        invariant to scale, topology, and language but degrades under compositional OOD (multiple
        simultaneous conflicts), consistent in direction with the weakest closed-loop bucket.
\end{itemize}
% Positioning: the earlier 540-case MVT is introduced only as an underpowered diagnostic that exposed
% nondiscriminative substrate behavior, NOT as STAG evidence (no superiority claim from it).
